\documentclass[guide]{pol-ren}

\title[Pronunciation]{Singing Polish and Polish Latin}
\subtitle{A pronunciation guide to the editions}
\shortcomposer{Polish Early Music}
\forces{Modern Polish · Polish of the sixteenth century · Old Polish · Latin as sung in Poland}
\edition{First edition}{2026}
\slug{pronunciation}
\repository{github.com/mderdun/pol-ren}

\begin{document}
\maketitlepage
\colophon

\section*{How to use this guide}

\initial{T}{he editions} print each text in its own spelling and leave the
sounds to this guide. An edition’s notes mention only what is peculiar to its
text. Sounds are given in the International Phonetic Alphabet between square
brackets, with a stress mark before the stressed syllable and a full stop
between syllables where the division matters for underlay.

Each section gives the series default first: what the editions assume and what
their underlay is built on. Historical options follow, with the evidence for
each. Take an option only if the whole choir can keep it, and say so in the
programme. A mixed pronunciation is worse than either.

\section*{Modern Polish}

Polish spelling is regular. A letter or a pair of letters has one sound,
with a few rules of context given below. Every vowel is a full vowel; nothing
is reduced to a neutral \ipa{ə}. Stress falls on the next-to-last syllable.

\begin{soundtable}
\tablegroup{vowels}
a & \ipa{a} & \pl{matka} \ipa{ˈmat.ka}\\
e & \ipa{ɛ}, open & \pl{dzień} \ipa{dʑɛɲ}\\
i & \ipa{i} & \pl{miły} \ipa{ˈmi.wɨ}\\
o & \ipa{ɔ}, open & \pl{noc} \ipa{nɔts}\\
ó, u & \ipa{u} & \pl{Bóg} \ipa{buk}, \pl{używać} \ipa{uˈʐɨ.vatɕ}\\
y & \ipa{ɨ}, an i with the tongue drawn back & \pl{był} \ipa{bɨw}\\
ą & nasal o: \ipa{ɔw̃} (see below) & \pl{są} \ipa{sɔw̃}\\
ę & nasal e: \ipa{ɛw̃} (see below) & \pl{męka} \ipa{ˈmɛŋ.ka}\\[0.4em]
\tablegroup{consonants}
c & \ipa{ts} & \pl{noc} \ipa{nɔts}\\
ch, h & \ipa{x}, as in \textit{loch} & \pl{chytrość} \ipa{ˈxɨ.trɔɕtɕ}\\
cz & \ipa{tʂ}, a hard \textit{ch} & \pl{czart} \ipa{tʂart}\\
ć, ci + vowel & \ipa{tɕ}, a soft \textit{ch} & \pl{ciemność} \ipa{ˈtɕɛm.nɔɕtɕ}\\
dz & \ipa{dz} & \pl{dzwon} \ipa{dzvɔn}\\
dź, dzi + vowel & \ipa{dʑ}, a soft \textit{j} & \pl{chodzi} \ipa{ˈxɔ.dʑi}\\
dż & \ipa{dʐ}, a hard \textit{j} & \pl{dżdżysty} \ipa{ˈdʐdʐɨs.tɨ}\\
g & \ipa{ɡ}, always hard & \pl{Boga} \ipa{ˈbɔ.ɡa}\\
j & \ipa{j} & \pl{jesteś} \ipa{ˈjɛs.tɛɕ}\\
ł & \ipa{w} & \pl{był} \ipa{bɨw}\\
ń, ni + vowel & \ipa{ɲ} & \pl{Panie} \ipa{ˈpa.ɲɛ}\\
r & \ipa{r}, rolled & \pl{Kryste} \ipa{ˈkrɨs.tɛ}\\
rz, ż & \ipa{ʐ}, as in \textit{measure} with the tongue drawn back & \pl{strażem} \ipa{ˈstra.ʐɛm}\\
s & \ipa{s} & \pl{syn} \ipa{sɨn}\\
ś, si + vowel & \ipa{ɕ}, a soft \textit{sh} & \pl{się} \ipa{ɕɛ}\\
sz & \ipa{ʂ}, a hard \textit{sh} & \pl{nasze} \ipa{ˈna.ʂɛ}\\
w & \ipa{v} & \pl{wieczny} \ipa{ˈvʲɛtʂ.nɨ}\\
z & \ipa{z} & \pl{złych} \ipa{zwɨx}\\
ź, zi + vowel & \ipa{ʑ}, a soft \textit{zh} & \pl{zima} \ipa{ˈʑi.ma}\\
\end{soundtable}

\subsection{Soft consonants and the silent i} Between a consonant and a
vowel, \textit{i} only softens the consonant. It is not a vowel and never
takes a note: \pl{się} is one syllable \ipa{ɕɛ}, \pl{Panie} two
\ipa{ˈpa.ɲɛ}, \pl{miłosierny} four \ipa{mi.wɔˈɕɛr.nɨ}. Elsewhere \textit{i} is
a vowel. Latin and borrowed words are different (see the Latin section).

\subsection{Nasal vowels} \pl{Ą} and \pl{ę} are true nasal vowels before
\textit{s, z, sz, ż, ś, ź, f, w, ch} and at the end of a word: an oral vowel
followed by a nasal glide, \ipa{ɔw̃} and \ipa{ɛw̃}. Before a stop they become
the vowel and a nasal consonant: \pl{męka} \ipa{ˈmɛŋ.ka}, \pl{święty}
\ipa{ˈɕfʲɛn.tɨ}, \pl{ząb} \ipa{zɔmp}. Before \textit{l} and \textit{ł} they
lose their nasality. At the end of a word \pl{ę} is usually plain \ipa{ɛ} in
speech. On a long note, hold the oral vowel and let the nasal glide come at
the end of the note, as it does in speech.

\subsection{Voicing} A voiced consonant at the end of a word is voiceless:
\pl{Bóg} \ipa{buk}. In a cluster the consonants agree with the last one:
\pl{wszytko} \ipa{ˈfʂɨt.kɔ}, \pl{swej} \ipa{sfɛj}, \pl{przez} \ipa{pʂɛs}.
A one-letter preposition joins the next word: \pl{w ciemności}
\ipa{fˈtɕɛm.nɔɕ.tɕi}, \pl{z łaski} \ipa{ˈzwas.ki}.

\subsection{Stress} On the next-to-last syllable: \pl{poprośmy}
\ipa{pɔˈprɔɕ.mɨ}, \pl{używają} \ipa{u.ʐɨˈva.jɔw̃}. Short unstressed words
(\pl{się, mi, cię, go, nas}) lean on the word before.

\section*{Polish of the sixteenth century}

\subsection{Default} Modern pronunciation of the old forms as printed. The
editions keep forms such as \pl{zmierzka, naszem, swoję, wszytki, jenż}
because they change rhyme, syllable count or sense. Their sounds changed
gradually and unevenly over two centuries, and singing them in modern
pronunciation is what Polish choirs do. Read \pl{naszem} as \ipa{ˈna.ʂɛm},
\pl{swoję} as \ipa{ˈsfɔ.jɛ}.

The following historical options are attested for the Polish of about 1550.
They are listed from the best attested and easiest to the least certain.

\subsection{Ł as a dark l} In the sixteenth century \textit{ł} was a lateral,
the dark \ipa{ɫ} of English \textit{full}. Its change to \ipa{w} began around
1600 and became standard only in the twentieth century; stage and educated
speech kept \ipa{ɫ} until then. This is the safest historical option:
\pl{był} \ipa{bɨɫ}, \pl{miły} \ipa{ˈmi.ɫɨ}.

\subsection{Rz as a trilled fricative} Written \textit{rz} began as a soft
\textit{r}, and by the fourteenth century it had become a trill with friction,
\ipa{r̝}, the sound of Czech \textit{ř} in \textit{Dvořák}. It merged with
\textit{ż} gradually between the sixteenth and eighteenth centuries, earlier
in some places than in others. For Kraków around 1550, \ipa{r̝} is
defensible: \pl{rzeczy} \ipa{ˈr̝ɛ.tʂɨ}, \pl{przez} \ipa{pr̝̊ɛs}.

\subsection{Stress over the stress group} Stress already fell on the
next-to-last syllable, but it was counted over the whole group, including a
short word leaning on the one before. \pl{Stało się} was stressed
\ipa{staˈɫɔ ɕɛ} where modern Polish has \ipa{ˈsta.wɔ ɕɛ}. This matters at a
cadence where \pl{się, mi} or \pl{cię} follows a two-syllable word.

\subsection{Nasal vowels} In the sixteenth century \pl{ą} was still a nasal
a, close to \ipa{ɒ̃}, and reached its modern nasal o only in the seventeenth
and eighteenth centuries. \pl{Ę} was more open than now, \ipa{æ̃} to
\ipa{ɛ̃}.

\subsection{Narrowed vowels} The long vowels of Old Polish had become
narrowed vowels (\pl{samogłoski pochylone}) by about 1500: a narrowed a
between a and o, \ipa{ɒ}; a narrowed e between e and i or y; a narrowed o
between o and u, which became the modern \pl{ó} \ipa{u}. Which words had
them is a fact about each word, and the prints are no reliable guide.
Printers usually marked the narrowed \pl{é} and \pl{ó} with an accent, but for
a they usually did the opposite, accenting the clear vowel and leaving the
narrowed one plain (Statorius, 1568: \pl{bábá, spáć} against \pl{pan, sad}).
Practice was inconsistent. A choir that wants narrowed vowels needs a list
for the text in hand; the editions do not mark them.

\subsection{Reading the old prints} The diplomatic texts in the critical
editions keep the printers’ spelling. These are the differences that matter:

\begin{soundtable}
y before a vowel, final y & \textit{j} & \pl{Yeſu} = \pl{Jesu}, \pl{pokoy} = \pl{pokój}\\
y, i after p, b, m, w & softness, as modern \textit{i} & \pl{zmyerzka} = \pl{zmierzka}, \pl{ſwyętego} = \pl{świętego}\\
v at the start of a word & \textit{u} & \pl{Vżywáyą} = \pl{używają}\\
ſ & \textit{s} & \pl{ſwey} = \pl{swej}\\
ſſ & \textit{sz} & \pl{náſſym} = \pl{naszym}\\
á & usually the clear a & \pl{Bogá} = \pl{Boga}\\
ē, ā & a vowel followed by m or n & \pl{zduchē} = \pl{z Duchem}\\
words run together & a preposition joined to its noun & \pl{wćyemnośći} = \pl{w ciemności}\\
\end{soundtable}

\section*{Old Polish}

\pl{Bogurodzica} is sung in Woronczak’s normalised text in modern
pronunciation. That is the usual practice, and the edition’s underlay is
built on it. The Polish of about 1400 differed a good deal. Stress fell on
the first syllable and moved to the next-to-last only around 1500. Vowels
were long or short. There was a single nasal vowel, written \pl{ǫ} in the
diplomatic text, which later split into \pl{ą} and \pl{ę}. \textit{Ł} was
\ipa{ɫ} and \textit{rz} was \ipa{r̝}. A sung reconstruction could be built
from the historical grammars, but no performing tradition has tested one,
and the editions do not print one.

\pl{Kyrie eleison} is Greek: \textit{y} is \ipa{i}, and the refrain has seven
syllables, \ipa{ˈki.ri.ɛ ɛˈlɛ.i.sɔn}.

\section*{Latin as sung in Poland}

\subsection{Default} The Polish tradition. Polish choirs sang Latin with
Polish sounds for some letters, and Polish churches and schools kept many of
them. It differs from the Roman pronunciation that most choirs know in about
a dozen spellings. Everything not listed is as in Roman Latin.

\begin{latintable}
c before e, i, y, ae, oe & \ipa{ts} & \ipa{tʃ} & \la{caelum} \ipa{ˈtsɛ.lum}, \la{principio} \ipa{prinˈtsi.pi.ɔ}\\
sc before e, i & \ipa{sts} & \ipa{ʃ} & \la{scio} \ipa{ˈstsi.ɔ}\\
ch & \ipa{x} & \ipa{k} & \la{Rachel} \ipa{ˈra.xɛl}\\
h & \ipa{x} & silent & \la{Herodis} \ipa{xɛˈrɔ.dis}\\
g & \ipa{ɡ}, always & \ipa{dʒ} before e, i & \la{angelum} \ipa{ˈaŋ.ɡɛ.lum}\\
gn & \ipa{ɡn} & \ipa{ɲ} & \la{agnus} \ipa{ˈaɡ.nus}\\
qu & \ipa{kv} & \ipa{kw} & \la{qui} \ipa{kvi}\\
ngu before a vowel & \ipa{ŋɡv} & \ipa{ŋɡw} & \la{sanguis} \ipa{ˈsaŋ.ɡvis}\\
s between vowels & \ipa{z} & \ipa{s} or \ipa{z} & \la{misit} \ipa{ˈmi.zit}, \la{Iesu} \ipa{ˈjɛ.zu}\\
ti before a vowel & \ipa{tsi} & \ipa{tsi} & \la{gratia} \ipa{ˈɡra.tsi.a}\\
z & \ipa{z} & \ipa{dz} & \la{zelus} \ipa{ˈzɛ.lus}\\
ae, oe & \ipa{ɛ} & \ipa{ɛ} & \la{saecula} \ipa{ˈsɛ.ku.la}\\
y & \ipa{i} & \ipa{i} & \la{Kyrie} \ipa{ˈki.ri.ɛ}\\
\end{latintable}

\subsection{Vowels and syllables} Five pure vowels, none reduced. Unlike
Polish, \textit{i} before a vowel is a syllable of its own:
\la{gratia} has three, \la{Maria} three, \la{filios} three. The
sixteenth-century Polish forms of these words show it: \pl{gracyja, Maryja}.

\subsection{Ci, si, ti, di, ri before a vowel} These stay hard, \ipa{tsi},
\ipa{si}, \ipa{ti}, \ipa{di}, \ipa{ri}, and are not softened to the Polish
\ipa{tɕ}, \ipa{ɕ}. The Polish loans of the period have hard consonants
(\pl{lekcyja, pasyja, historyja}). \textit{N} before \textit{i} is another
matter. Polish speakers soften it, and the loans have a soft n
(\pl{unija, Antoni}), so it may stay soft.

\subsection{Ti} \ipa{tsi} with a full \ipa{i}: \la{exspectatione}
\ipa{ɛk.spɛk.ta.tsiˈɔ.nɛ}. Not after \textit{s, t, x}: \la{hostia}
\ipa{ˈxɔs.ti.a}. Polish schools now often teach \ipa{ti}; the church
tradition and the period have \ipa{ts}.

\subsection{S between vowels} \ipa{z}. The Polish loans voice it
(\pl{muzyka, filozof, Jezus}), and the series follows them.

\subsection{G before e and i} Hard \ipa{ɡ}. Medieval Polish readers said
\ipa{j} here, which is why \la{angelus} became \pl{anioł} and
\la{evangelium} \pl{ewanjelija}; the hard \ipa{ɡ} is the later learned norm.
When it took over in sung Latin is not known. For fifteenth-century
repertory, \ipa{j} is a historical option.

\subsection{Qu} \ipa{kv}, which Polish voicing rules turn into \ipa{kf} in
quick speech. Either is right.

\subsection{Stress} By the rules of Latin, not of Polish: the next-to-last
syllable if it is long, otherwise the one before. Jerzy Liban’s treatise on
church accentuation (Kraków, c.\,1539) teaches these rules, and Polish kept
Latin stress even in borrowed words (\pl{muzyka} \ipa{ˈmu.zɨ.ka}). Do not move
the stress to the next-to-last syllable: \la{Dóminus, ángelum, erípuit,
fílios, nóluit}.

\section*{Further reading}

\begin{literature}
\lit Długosz-Kurczabowa, Krystyna, and Stanisław Dubisz. \pl{Gramatyka
  historyczna języka polskiego}. Warsaw: Wydawnictwa Uniwersytetu
  Warszawskiego, 2006.
\lit German, Jan. \pl{Drogi przenikania zapożyczeń greckich i łacińskich do
  języka staropolskiego}. PhD dissertation, Uniwersytet Jagielloński, 2023.
\lit Klemensiewicz, Zenon. \pl{Historia języka polskiego}. Warsaw: PWN, 1999.
\lit Kuraszkiewicz, Władysław. \pl{Gramatyka historyczna języka polskiego}.
  Warsaw: PZWS, 1972.
\lit Liban, Jerzy. \la{De accentuum ecclesiasticorum exquisita ratione}.
  Kraków: Scharffenberg, c.\,1539.
\lit Statorius-Stojeński, Piotr. \la{Polonicae grammatices institutio}.
  Kraków, 1568.
\lit Urbańczyk, Stanisław. ‘Z zagadnień staropolskich, 4: Ańjoł, jenerał
  itp.’ \textit{Język Polski} 32 (1952): 127–29.
\lit Walczak, Bogdan. \pl{Zarys dziejów języka polskiego}. Wrocław:
  Wydawnictwo Uniwersytetu Wrocławskiego, 1999.
\lit Woronczak, Jerzy, Ewa Ostrowska and Hieronim Feicht. \pl{Bogurodzica}.
  Wrocław: Ossolineum, 1962.
\end{literature}

\end{document}
